% Fragmento público generado desde propietarios íntegros.
% Residencia: 52.2.3.
% Fuente: ../incoming/radion_monografia_20260827/manuscrito/sections/03_cierre_exacto.tex:104-253
\subsubsection{Transition cocycle between the direct and torsional sections}

\paragraph{definition operatoria of the carry \(\varepsilon_R\)}

\begin{definition}[Real carry of the radion]
The unitary output of the radion is the transition cocycle
\begin{equation}
\epsone:=S_T-S_E
=\Phi_T-D_E
=\frac{20}{81}\Delta_4-operatorname{rem}_1(S_E).
\label{eq:epsilon-def}
\end{equation}
\end{definition}

The equation contains a subtraction, but it is not a fit to the target. Its two operands \(S_T\) and \(S_E\) are generated separately. The independence demonstrated is independence from \(180/\pi\) and from a target \(\varepsilon\); algebraic independence between a difference and its summands is not asserted.

\begin{theorem}[Exact unitary closure]
With \(S_E\), \(\Phi_T\), and \(\epsone\) defined by
\eqref{eq:seccion-electrica}, \eqref{eq:delta4}, \eqref{eq:rho-twoway}, and
\eqref{eq:epsilon-def}, one has exactly
\begin{equation}
\boxed{1=S_E-\Phi_T+\epsone.}
\label{eq:cierre-unitario}
\end{equation}
\end{theorem}
\begin{proof}
By \eqref{eq:SE-intervalo}, \(D_E=S_E-1\). By definition,
\(\epsone=\Phi_T-D_E\). Then
\[
S_E-\Phi_T+\epsone
=S_E-\Phi_T+\Phi_T-(S_E-1)=1.
\]
The strength of the proof lies not in this elementary cancellation, but in the fact that the
prior genealogy fixes both operands without using the closure as input.
\end{proof}

\paragraph{Publication of the closure on the angular chart}

The full-turn chart multiplies a unit fraction by \(360/T_+=180/\pih\). In particular,
\[
\epsR^\circ=\frac{360^\circ}{T_+}\epsone.
\]
Multiplying \eqref{eq:cierre-unitario} by \(360^\circ/T_+\) and using
\[
\frac{360}{T_+}S_E
=360\left(\frac5{36}+\frac{25}{9}\alphah\right)
=50+1000\alphah=50+A
\]
produce the angular closure.

\begin{theorem}[Exact equation of the radion]
\BigRadionEquation
\end{theorem}

The equation says something more precise than \enquote{the radian is
\(57.2957\ldots{}^\circ\)}. It decomposes that publication into four operations
with a reference proof:
\begin{center}
\begin{tabularx}{.96\textwidth}{>{\bfseries}l X}
\toprule
\(50^\circ\) & second phase of the wheel; critical hinge on a declared chart;\\
\(A^\circ\) & common coordinate of the parangular publication;\\
\(-\frac{20}{81}\Delta_4^\circ\) & orientation-transport of global torsion;\\
\(+\epsR^\circ\) & real carry that mates the two lifts.\\
\bottomrule
\end{tabularx}
\end{center}

\paragraph{Convergence by depth of refinement}

At depth \(N\), the generator delivers rational cylinders
\(p_N\to\pih\) and \(e_N\to e_{\HMT}\). Define
\[
S_{E,N}=2p_N\left(\frac5{36}+\frac{25}{9}\alphah\right),
\]
\[
\Delta_{4,N}=\frac{p_N-e_N\log p_N}{270}
\left(1+\frac{p_N}{729}\right),\qquad
S_{T,N}=1+\frac{20}{81}\Delta_{4,N},
\]
and
\[
\mathcal E_N=S_{T,N}-S_{E,N}.
\]
Continuity of the functional on \(3<p<4\), \(2<e<3\), together with the
monotone derivatives, transports each pair of cylinders to a certified
interval. At a depth of \(45\) blocks, its width is less than
\(4.81\times10^{-130}\).

Each nine-level return contains six trits per level. Therefore, the refinement rate is
\begin{equation}
729^{-9}=3^{-54}.
\label{eq:contraccion-54}
\end{equation}
This quantity controls the width of the intervals; it is not the value of
\(\epsR\).

The subtraction of triads stabilizes
\begin{center}
\ttfamily
000|000|000|896|802|822|044|562|981|999|050|695|020|651|286|413
\end{center}
and the carry prefix
\begin{center}
\ttfamily
0,0,0,0,0,-1,0,-1,-1,-1,0,-1,0.
\end{center}
This is the concrete mechanism of decimal memory: not a chosen tail, but
the recurrence \eqref{eq:subacarreo} applied to previously
constructed operands.

\paragraph{Evaluation of the canonical values}

The coinductive output is
\begin{align}
\epsone
&=8.9680282204456298199905069502065128641372736205009\ldots
\times10^{-10},\\
\epsR^\circ
&=5.1383016758575282071685164622645947489908574400054\ldots
\times10^{-8}\,{}^\circ.
\label{eq:epsilon-valores}
\end{align}
The complete publication results
\[
\frac{180^\circ}{\pih}
=57.2957795130823208767981548141051703324054724665643\ldots{}^\circ.
\]

\begin{exactbox}
\begin{align*}
50^\circ+A^\circ
&=57.29735256928380099728510547238066299956\ldots{}^\circ,\\
\frac{20}{81}\Delta_4^\circ
&=0.00157310758449687906223272996065728980465\ldots{}^\circ,\\
50^\circ+A^\circ-\frac{20}{81}\Delta_4^\circ
&=57.29577946169930411822287274242000570976\ldots{}^\circ,\\
\epsR^\circ
&=0.00000005138301675857528207168516462264594749\ldots{}^\circ,\\
\text{final sum}
&=57.29577951308232087679815481410517033241\ldots{}^\circ.
\end{align*}
\end{exactbox}

